\documentclass[10pt,a4paper]{amsart}
\usepackage[margin=19mm]{geometry}
\usepackage{amsmath,amssymb,mathtools}
\usepackage[scr=rsfs,cal=euler]{mathalfa}
\usepackage{xfrac}
\usepackage[unicode,hidelinks,bookmarks=false]{hyperref}
\newcommand{\ie}{i.e.\ }
\newcommand{\cf}{cf.\ }
\newcommand{\resp}{respectively\ }
\newcommand{\Subsection}{\subsection}
\providecommand{\nobreakdash}{\nobreak\hbox{-}}
\setlength{\emergencystretch}{2em}
\multlinegap0pt
\usepackage{amssymb}
\usepackage{mathtools}
\newtheorem{thm}{Theorem}
\def\as{\coloneqq}
\def\col{\colon}
\def\eps{\varepsilon}
\renewcommand{\hat}{\widehat}
\newcommand{\hlf}[1]{\frac{#1}{2}}
\def\l{\left}
\def\nat{\mathbb{N}}
\def\ol{\overline}
\def\on{\overline{n}}
\def\r{\right}
\newcommand{\rec}[1]{\frac{1}{#1}} 
\def\rls{\mathbb{R}}
\def\un{\underline{n}}
\title{The Kaplan--Meier estimator as a limit function of Efron's self-consistency iterations}
\author{Miroslav Ba\v{c}\'ak}
\date{Source: arXiv:2607.27068v2 (CC BY 4.0)}
\begin{document}
\maketitle

\noindent\emph{Source and licence.} The Kaplan--Meier estimator as a limit function of Efron's self-consistency iterations. Version arXiv:2607.27068v2 (CC BY 4.0), distributed by arXiv under the Creative Commons Attribution 4.0 International licence, http://creativecommons.org/licenses/by/4.0/, as recorded in the arXiv licence metadata for this exact version and shown on the abstract page https://arxiv.org/abs/2607.27068v2. Full licence notice: \texttt{licenses/arxiv-2607.27068-CC-BY-4.0.txt}.\par\smallskip

\noindent\textit{Editorial scope.} Counted unit: the article's thm thm:main (source0.tex lines 200--201). The article's complete original proof of it is delivered with it (source0.tex lines 202--289, one \texttt{proof} environment of 88 lines). No mathematics is written, repaired or re-derived: the statement and the proof are the article's own lines, in the article's own notation, and no retained line is substituted or re-flowed.  Retained uncounted as the source-local context the proof needs: lines 129--131; lines 139--141.  Nothing was omitted from any retained span, so the package carries no omission record.  No retained line contains a citation, so the wrapper carries no bibliography and no bibliography entry.  No retained line contains a figure environment, an image-inclusion command or drawing code, so no image is delivered and the package declares no asset dependency.  Editorial formatting only, in addition: an AMS \texttt{amsart} preamble that restates the article's own declarations and macros used by the retained lines, namely the environments \texttt{thm} on separate counters, together with the standard packages those lines need.  The retained environments print in this document as Theorem 1 in order of appearance, whereas the article prints the same items with the numbering of its own full document.  The complete native source of the article as distributed in the arXiv e-print for this exact version is bundled unchanged as \texttt{sources/206318/source0.tex}.
\begin{equation} \label{eq:self}
 \hat{S}(t) = e(t) + \rec{m} \sum_{\substack{t_i\le t \\ \delta_i=0}}\frac{\hat{S}(t)}{\hat{S}\l(t_i\r)},\qquad t\in\rls,
\end{equation}

\begin{equation} \label{eq:iter}
 \hat{S}_{k+1}(t) \as e(t) + \rec{m} \sum_{\substack{t_i\le t \\ \delta_i=0}}\frac{\hat{S}_k(t)}{\hat{S}_k\l(t_i\r)},\qquad t\in\rls,\;k\in\nat,
\end{equation}

\begin{thm}[Convergence of Efron's self-consistency iterations] \label{thm:main} Given an arbitrary survival function~$\hat{S}_1$ as a starting point of the iterative algorithm in~\eqref{eq:iter}, the sequence $\l(\hat{S}_k\r)$ converges uniformly on $\l(-\infty,t_m\r)$ to a~unique self-consistent estimator, which coincides with the Kaplan--Meier estimator.
\end{thm}

\begin{proof}
Given $n\in\{1,\dots, m\},$ denote
\begin{align*}
 \un & \as \min\l\{i\in\{1,\dots,m\}\col t_i = t_n  \r\} ,\\
 \overline{n} & \as \max\l\{i\in\{1,\dots,m\}\col t_i = t_n  \r\} .
\end{align*}
For the reader's convenience, we divide the proof into several steps. 

\textbf{Step 1:} Observe that $\hat{S}_k$ is a survival function for each $k\in\nat.$ We proceed by induction. The starting function~$\hat{S}_1$ is a survival function and assume that $\hat{S}_k$ is a survival function for some fixed $k\in\nat.$ Then~$\hat{S}_{k+1}$ is a survival function, too. Indeed, it is easy to see that~$\hat{S}_{k+1}(t)=1$ for $t\in\l(-\infty,t_1\r),$ and that $\hat{S}_{k+1}$ is a non-increasing right-continuous function.

\textbf{Step 2:} Let us now prove that the sequence $\l(\hat{S}_k(t)\r)$ converges, as $k\to\infty,$ for each $t\in\l(-\infty,t_m\r).$ If $t\in\l(-\infty,t_1\r),$ then $\hat{S}_k(t)=e(t)=1$ for each $k\in\nat.$ If $m=1,$ we are done. Let also $m>1$ and assume for some $n\in\{1,\dots,m-1\}$ the sequence $\l(\hat{S}_k(t)\r)$ converges for each $t\in\l(-\infty,t_n\r).$ Without loss of generality we may and do assume $n=\ol{n}.$ 

\textbf{Step 2a:} We will show that the sequence $\l(\hat{S}_k\l(t_n\r)\r)$ converges, as $k\to\infty.$ To this end observe that each new iteration $\hat{S}_{k+1}\l(t_n\r)$ depends only on the values of $\hat{S}_k$ at $t_1,\dots, t_n.$ Set therefore
\begin{equation} \label{eq:bk-notation}
b_k\as \rec{m} \sum_{\substack{t_i< t_n \\ \delta_i=0}}\rec{\hat{S}_k\l(t_i\r)},\qquad\text{for each } k\in\nat.                                                                                                                                                                                                                                                                                                                                                                                                                                                                \end{equation}
Here we use the convention that an ``empty sum'' is zero, that is, if $\delta_i=1$ for every $i$ with $t_i< t_n,$ then $b_k=0.$ With this notation the iterative scheme~\eqref{eq:iter} reads
\begin{equation} \label{eq:iter-bk}
 \hat{S}_{k+1} \l(t_n\r) = e\l(t_n\r) +\frac{\zeta_n}{m} + b_k \hat{S}_k \l(t_n\r), 
\end{equation}
for each $k\in\nat.$ Set $a\as e\l(t_n\r)+\frac{\zeta_n}{m}=\frac{m-n}{m}+\frac{\zeta_n}{m}.$ Observe that $a>0,$ since $n<m.$ Finally, set $s_k \as \hat{S}_k\l(t_n\r).$ With this notation at hand, one can write \eqref{eq:iter-bk} compactly as
\begin{equation*}
 s_{k+1}= a + b_k s_k.
\end{equation*}

By the assumption above, the sequence $\l(b_k\r)$ converges to some $b\in(0,\infty),$ as $k\to\infty.$ If $b>1,$ then
\begin{equation*}
 s_{k+1}-s_k=a+ \l(b_k-1\r)s_k \ge a,\qquad \text{for } k \text{ big enough,}
\end{equation*}
which contradicts the boundedness of $\l(s_k\r).$ Likewise, if $b=1,$ then
\begin{equation} \label{eq:adhoc}
 s_{k+1}-s_k=a+ \l(b_k-1\r)s_k \to a,\qquad \text{as } k\to\infty,
\end{equation}
since $0\le s_k\le 1,$ for every $k\in\nat.$ But \eqref{eq:adhoc} actually contradicts again the boundedness of the sequence~$\l(s_k\r).$ 

We therefore conclude $b<1$ and set $c\as\frac{a}{1-b}>0.$ We will show that the sequence $\l(s_k\r)$ converges to~$c.$ Given $\eps>0,$ we want to find a $p_0\in\nat$ such that $\l|s_p-c \r|<\eps$ for every $p\ge p_0.$ To this end, consider
\begin{equation*}
 s_{k+1}-c = a + b_k s_k - c = c(1-b) + b_k s_k - c = b_k\l(s_k-c\r) + \l(b_k-b\r)c, 
 \end{equation*}
 and estimate
 \begin{equation} \label{eq:estim1}
 \l|s_{k+1}-c \r| \le b_k \l|s_k-c\r| + \l|b_k-b\r|c, 
\end{equation}
for every $k \in\nat.$ Since the sequence $\l(b_k\r)$ converges to $b<1,$ there exist a constant $r\in (b,1)$ and an index $k_0\in\nat$ such that simultaneously 
\begin{equation*}
b_k<r,\qquad r^k(c+1)<\hlf{\eps},\qquad \l|b_k-b\r|c<\hlf{\eps}(1-r),
\end{equation*}
for every $k\ge k_0.$ Then we claim that $p_0\as 2k_0$ is the desired index. Indeed, let $p\ge p_0.$ Then there exist $k\ge k_0$ and $l\ge k_0$ such that $p=k+l.$ By iterating the estimate~\eqref{eq:estim1}, we obtain
\begin{align*}
 \l|s_{k+1}-c \r| & \le r \l|s_k-c\r| + \l|b_k-b\r|c \\ \l|s_{k+2}-c \r| & \le r^2 \l|s_k-c\r| + r \l|b_k-b\r|c + \l|b_{k+1}-b\r|c \\ \vdots \\
 \l|s_{k+l}-c \r| & \le r^l \l|s_k-c\r| + \sum_{j=0}^{l-1} r^{l-1-j}\l|b_{k+j}-b\r|c.
\end{align*}

Hence, using the fact that $s_k \in [0,1],$ we get
\begin{equation*}
 \l|s_p -c \r| = \l|s_{k+l}-c \r| \le r^l (c+1) +\sum_{j=0}^{l-1} r^{l-1-j}\hlf{\eps}(1-r) \le\hlf{\eps} + \hlf{\eps}(1-r) \sum_{j=0}^{l-1} r^j \le\hlf{\eps} + \hlf{\eps}(1-r) \sum_{j=0}^\infty r^j =\eps,
\end{equation*}
and this shows that the sequence $\l(\hat{S}_k\l(t_n\r)\r)$ converges, as $k\to\infty.$

\textbf{Step 2b:} Next we show that the sequence $\l(\hat{S}_k(t)\r)$ converges for each $t\in\l(t_n,t_{n+1}\r).$ Note that the empirical survival function $e$ is obviously constant on this interval, namely $e(t)=\frac{m-n}{m},$ for every $t\in\l(t_n,t_{n+1}\r).$ We keep using the notation~\eqref{eq:bk-notation}. By the definition of the iterative algorithm~\eqref{eq:iter} we have
\begin{equation*} 
 \hat{S}_{k+1}(t) = e(t) + \l(\frac{\zeta_n}{m}\rec{\hat{S}_k\l(t_n\r)}+ \rec{m} \sum_{\substack{t_i< t_n \\ \delta_i=0}}\rec{\hat{S}_k\l(t_i\r)} \r) \hat{S}_k(t) = \frac{m-n}{m} + \l(\frac{\zeta_n}{m}\rec{\hat{S}_k\l(t_n\r)}+ b_k \r) \hat{S}_k(t) .
\end{equation*}
We already know by Step 2a that
\begin{equation*}
 \frac{\zeta_n}{m}\rec{\hat{S}_k\l(t_n\r)}+ b_k \to\frac{\zeta_n}{m} \frac{1-b}{\frac{m-n+\zeta_n}{m}}+b,\qquad \text{as } k\to\infty.
\end{equation*}
One can therefore repeat Step 2a again with $t$ in place of $t_n$ and conclude that $\l(\hat{S}_k(t)\r)$ converges to
\begin{equation} \label{eq:adhoc-limit2}
 \frac{m-n}{m}\rec{1-\l(\frac{(1-b)\zeta_n}{m-n+\zeta_n}+b\r)} = \frac{m-n+\zeta_n}{m}\rec{1-b}.
\end{equation}
In summary, assuming that the sequence $\l(\hat{S}_k\r)$ converges on $\l(-\infty,t_n\r),$ we got its convergence $\l[t_n,t_{n+1}\r).$

\textbf{Step 3:} We can proceed inductively and prove the convergence of the sequence $\l(\hat{S}_k\r)$ on $\l(-\infty,t_m\r).$ 

By inspecting the above arguments, one can see that the convergence is even \emph{uniform} in $t\in\l(-\infty,t_m\r),$ since on each of the intervals $\l(t_n,t_{n+1}\r)$ the convergence was independent of~$t.$

\textbf{Step 4:}  It follows immediately that the limit function, which we denote by $\hat{S},$ is a (piece-wise constant) non-increasing function. We now show that it is right-continuous. To this end, consider an index $n\in\{1,\dots, m-1\},$ without loss of generality $n=\on,$ and compare the value $\hat{S}\l(t_n\r)$ with $\hat{S}\l(t\r),$ where $t\in \l(t_n,t_{n+1}\r).$ Indeed, $\hat{S}\l(t_n\r)=\frac{m-n+\zeta_n}{m}\rec{1-b}$ by Step 2a and using~\eqref{eq:adhoc-limit2} we arrive at $\hat{S}\l(t_n\r)=\hat{S}(t),$ for every $t\in \l(t_n,t_{n+1}\r),$ implying the function~$\hat{S}$ is right-continuous at~$t_n.$ It is straightforward that the limit function is a~fixed point of~\eqref{eq:self}, and we have hence established the existence of a self-consistent estimator.

\textbf{Step 5:} The uniqueness follows by the fact that the limit of the sequence $\l(\hat{S}_k(t)\r)$ is equal to~\eqref{eq:adhoc-limit2}, in particular, it is independent of the starting point $\hat{S}_1(t).$

\textbf{Step 6:} We show that the self-consistent estimator coincides with the Kaplan--Meier estimator. For $t\in\l(-\infty,t_1\r)$ it is obviously true, since $\hat{S}(t)=\hat{S}^\textrm{KM}(t)=1.$ Let $n\in\{1,\dots, m-1\}$ and without loss of generality assume $n=\on.$ Then we use recursion:
\begin{align*}
 \hat{S}\l(t_n\r) &  =\frac{m-n+\zeta_n}{m} \rec{1-\rec{m}\sum_{\substack{t_i< t_n \\ \delta_i=0}}\rec{\hat{S}\l(t_i\r)}} \\ & = \frac{m-n+\zeta_n}{m-n+\zeta_n+\eta_n} \frac{m-n+\zeta_n+\eta_n}{m} \rec{1-\rec{m}\sum_{\substack{t_i< t_n \\ \delta_i=0}}\rec{\hat{S}\l(t_i\r)}} \\  & = \l(1-\frac{\eta_n}{m-n+\zeta_n+\eta_n}\r)  \hat{S}\l(t_{\un-1}\r) \\ & = \l(1-\frac{\eta_n}{m-n+\zeta_n+\eta_n}\r) \hat{S}^\textrm{KM}\l(t_{\un-1}\r) \\ & = \hat{S}^\textrm{KM}\l(t_n\r).
\end{align*}
For $t\in\l(t_n,t_{n+1}\r),$ we have $\hat{S}(t)=\hat{S}\l(t_n\r)=\hat{S}^\textrm{KM}\l(t_n\r)=\hat{S}^\textrm{KM}(t),$ which shows that the self-consistent estimator $\hat{S}$ coincides with the Kaplan--Meier estimator $\hat{S}^\textrm{KM}$ on $\l(-\infty,t_m\r).$ 

This completes the proof.
\end{proof}

\end{document}
